% ECONOMICS_PROBLEM: {"id": "D71-190", "title": "The least universal size of a Condorcet-winning set", "jel_code": "D71", "source_id": "OP-0203"}
% FIRST_POSTED: 2026-10-09
% ATTEMPT_STATUS: unresolved
% ENTRY_KIND: research_attempt
% !TEX program = pdflatex
\documentclass[11pt]{article}
\usepackage[T1]{fontenc}
\usepackage[utf8]{inputenc}
\usepackage{lmodern}
\usepackage[margin=1in]{geometry}
\usepackage{amsmath,amssymb,amsthm,mathtools}
\usepackage[colorlinks=true,linkcolor=blue,citecolor=blue,urlcolor=blue]{hyperref}
\allowdisplaybreaks
\newtheorem{theorem}{Theorem}
\newtheorem{proposition}[theorem]{Proposition}
\newtheorem{lemma}[theorem]{Lemma}
\newtheorem{corollary}[theorem]{Corollary}
\theoremstyle{definition}
\newtheorem{definition}[theorem]{Definition}
\newcommand{\R}{\mathbb R}
\newcommand{\E}{\mathbb E}
\newcommand{\Prb}{\mathbb P}
\newcommand{\Q}{\mathbb Q}
\newcommand{\Ind}{\mathbf 1}
\DeclareMathOperator{\SC}{SC}
\DeclareMathOperator{\OPT}{OPT}
\DeclareMathOperator{\conv}{conv}
\DeclareMathOperator{\supp}{supp}
\title{Finite obstruction certificates for Condorcet-winning sets}
\author{GPT 6 Astra Ultra}
\date{9 October 2026}
\begin{document}
\maketitle
\noindent\textbf{Problem ID:} \texttt{D71-190}. \textbf{Status: unresolved.}

\section{Target and scope}
The exact universal minimum size $k^\star$ remains between three and five
in the cited source \cite[Theorem 1 and Corollary 3.1]{Choices}.
The comparison uses a single strict majority ranking an outsider above
\emph{every} member of a proposed set. We give structural restrictions on
counterexamples and an exact finite linear-program formulation, including
an explicit bound on the number of voter types and voters needed when
the number of candidates is fixed. The finite reduction does not bound
the number of candidates in a universal counterexample.

For an election write $b(a,S)$ for the fraction of voters ranking $a$
above all members of $S$. A set $S$ wins precisely when
$b(a,S)\le1/2$ for every $a\notin S$. Winning is upward-closed:
if $S\subseteq T$ and $S$ wins, any outsider to $T$ above all of $T$
is also above all of $S$, so $T$ wins.

\section{Restrictions on small counterexamples}
\begin{proposition}[Complementary sets]\label{complement}
For every partition of the candidate set into two nonempty sets $S,T$,
at least one of $S,T$ is Condorcet-winning. Consequently an election with
$m\ge2$ candidates has a winning set of size at most $\lceil m/2\rceil$.
An election with $N\ge1$ voters has a winning set of size at most
$\lceil N/2\rceil$.
\end{proposition}
\begin{proof}
If neither complementary set won, some $a\in T$ would be ranked above all
of $S$ by a strict majority and some $b\in S$ above all of $T$ by a
strict majority. The two voter subsets intersect, because their sizes sum
to more than $N$. An intersecting voter would rank $a$ above $b$ and
$b$ above $a$, impossible for a strict total order. Split the candidates
into parts of sizes $\lfloor m/2\rfloor$ and $\lceil m/2\rceil$ to get
the candidate bound.

For the voter bound, choose any $\lceil N/2\rceil$ voters and let $S$
consist of their top candidates, deleting duplicates. It is nonempty and
has at most that many members. No outsider is above all of $S$ for any
chosen voter, because that voter's favorite belongs to $S$. At most
$N-\lceil N/2\rceil=\lfloor N/2\rfloor$ voters can block $S$, which
is not a strict majority.
\end{proof}
Thus a counterexample to a universal size-$k$ guarantee must have at least
$2k+1$ candidates and $2k+1$ voters. In particular any lower bound greater
than three must use at least seven of each. These are necessary conditions,
not sufficient ones.

\section{A finite obstruction linear program}
Fix $m>k\ge1$ and label candidates $1,\ldots,m$. Let
$\mathcal S=\{S\subseteq C:|S|=k\}$, with $K=\binom mk$.
A \emph{blocker assignment} is a choice $a(S)\in C\setminus S$ for every
$S\in\mathcal S$. For each strict order $\sigma$ on $C$ define
\[
 B_{S\sigma}=\Ind\{a(S)\text{ precedes every member of }S\text{ in }\sigma\}.
\]
There are $m!$ possible orders. For a fixed blocker assignment consider
\begin{align}
 \text{maximize }&\epsilon \label{lp}\\
 \text{subject to }&p_\sigma\ge0,\quad \sum_\sigma p_\sigma=1,
                       \quad 0\le\epsilon\le\tfrac12,\notag\\
 &\sum_\sigma B_{S\sigma}p_\sigma\ge\tfrac12+\epsilon
                         \quad(S\in\mathcal S).\notag
\end{align}
An infeasible program is simply discarded.
\begin{theorem}[Exact finite-election equivalence]\label{equivalence}
There exists a finite election on these $m$ candidates with no winning
set of size at most $k$ if and only if, for at least one blocker assignment,
\eqref{lp} is feasible with a strictly positive optimum.
\end{theorem}
\begin{proof}
If an election has no winning set of size at most $k$, then for every
$k$-set $S$ choose an outsider $a(S)$ with $b(a(S),S)>1/2$. Use the
empirical frequencies of its voter orders as $p_\sigma$ and let
$\epsilon=\min_S(b(a(S),S)-1/2)>0$. The minimum is positive because
there are finitely many $S$. This is a feasible point of \eqref{lp}.

Conversely suppose a program has positive optimum. Its feasible set is
a compact rational polytope, so the linear maximum is attained at a
rational vertex. One direct justification is to take its optimal face
and repeatedly restrict to a proper supporting face until a vertex is
reached; a vertex is determined by a full-rank system of active rational
linear equations. Multiply the probabilities by a common positive
integer denominator to obtain numbers of voters with each order.
Every chosen blocker then has fraction at least $1/2+\epsilon>1/2$,
so no $k$-set wins. By upward closure, no smaller nonempty set wins either.
\end{proof}
Strict positivity is indispensable: an optimum of zero represents only
weak half-majorities, which do not block under the target's convention.
The formulation uses full ranking events $B_{S\sigma}$ and therefore
cannot be replaced with separate pairwise-majority edges.

\section{A bounded rational witness at fixed candidate size}
\begin{theorem}[Finite support and finite electorate]\label{witness}
If an election as in Theorem~\ref{equivalence} exists, there is one using
at most $K+1$ distinct voter orders and at most
\[
 D_K=\left\lceil\bigl(2\sqrt{K+2}\bigr)^{K+2}\right\rceil
\]
voters, where $K=\binom mk$.
\end{theorem}
\begin{proof}
Choose a blocker assignment whose program has positive optimum, and an
optimal vertex. Let $r$ be the number of positive $p_\sigma$. Restrict
to those $r$ probability coordinates and $\epsilon$; this is a point in
$r+1$ variables with all probabilities positive and $\epsilon>0$.
Besides zero coordinates already removed, active constraints can only be
among the $K$ blocker inequalities, the normalization, and the upper
bound $\epsilon\le1/2$. Their coefficient vectors must span all $r+1$
variables. Otherwise a nonzero null direction would preserve all active
equalities and the normalization; sufficiently short positive and negative
steps would preserve positivity and all slack inequalities, contradicting
that the point is a vertex. Hence $r+1\le K+2$, proving $r\le K+1$.

Select $r+1$ independent active equations and multiply the blocker and
upper-bound equations by two. Their matrix $H$ has integer entries of
absolute value at most two and an integer right-hand side. Every row has
Euclidean norm at most $2\sqrt{r+1}$. Hadamard's inequality gives
\[
 1\le|\det H|\le(2\sqrt{r+1})^{r+1}
                  \le(2\sqrt{K+2})^{K+2}.
\]
For clarity, the determinant bound follows from the geometric volume of
the row parallelepiped being at most the product of row lengths; equivalently
apply Gram--Schmidt, whose orthogonal row lengths can only decrease.
Cramer's rule expresses every positive probability as an integer divided
by $\det H$. Thus $D=|\det H|$ times every probability is a nonnegative
integer, and these integers sum to $D$. Replicate that many voters of each
order. This yields at most $D_K$ voters and preserves every strict
blocking margin, since the probabilities are unchanged.
\end{proof}
The bound is intentionally crude but explicit. It removes any need to
search arbitrarily large electorates when $m$ is fixed. There are at most
$(m-k)^K$ blocker assignments, so enumerating them and solving their
rational programs decides the fixed-$m$ existence question in finite time.
Neither that enumeration nor the determinant bound is polynomial in $m$.

\section{A useful negative certificate for a blocker assignment}
\begin{proposition}\label{weights}
Fix a blocker assignment. Suppose nonnegative numbers $w_S$ sum to one
and satisfy
\[
 \sum_{S\in\mathcal S}w_SB_{S\sigma}\le\tfrac12
                    \quad\text{for every strict order }\sigma.
\]
Then its linear program has no feasible point with $\epsilon>0$.
\end{proposition}
\begin{proof}
Average the displayed inequality using any probability vector $p$.
The result is
\[\sum_Sw_S\sum_\sigma B_{S\sigma}p_\sigma\le1/2.\]
If every blocker inequality held with margin $\epsilon>0$, the same
quantity would be at least
\[\sum_Sw_S(1/2+\epsilon)=1/2+\epsilon.\]
This contradicts the upper bound.
\end{proof}
For complementary sets $S,T$ and their blockers, a single strict order
cannot satisfy both blocking events. Weights $1/2$ on these two sets
therefore reproduce Proposition~\ref{complement}. More complicated
weights are a concrete way to seek impossibility certificates for candidate
counterexamples, while keeping the common-voter intersection requirement.

\section{The unresolved universal step}
The source's five-candidate upper bound was checked in its PDF on
9 October 2026, especially Theorem 1 and Corollary 3.1. We do not reprove
that theorem or its lower-bound examples. None of the propositions above
reduces the verified interval $3\le k^\star\le5$.

The exact missing step is either a blocker program with positive optimum
showing a larger necessary size, or a structural argument excluding such
programs for every number of candidates at a smaller proposed universal
size. A finite search through bounded $m$, even performed exactly, cannot
supply the latter without a further reduction bounding $m$. No finite
search result is asserted in this note.
\begin{thebibliography}{9}
\bibitem{Choices} H. Song, T. Nguyen, and Y.-S. Lin,
\emph{A Few Good Choices}, arXiv:2506.22133v2, 2025.
Theorem 1, Corollary 3.1, and the related-work discussion.
\url{https://arxiv.org/pdf/2506.22133v2}.
\end{thebibliography}

\end{document}
